% P15-U1: finite probability foundations, serving N-x-008.
\chapter{Finite Probability}
\label{sec:p15-finite-probability}

A finite probability model assigns weights to a specified set of outcomes.
From those weights one can compute probabilities of events, revise them after
an event of positive probability, and average numerical quantities. These
operations answer different questions: the probability of an observation
under a hypothesis is not the probability of that hypothesis after the
observation, and an average does not determine an individual outcome.

This chapter develops the finite part of the discrete treatment in
Grinstead and Snell's \emph{Introduction to Probability}, using the CHANCE
Project version dated 4 July 2006
\cite[secs.~1.2, 4.1 and 6.1]{GrinsteadSnellChance2006}.
The presentation fixes an underlying outcome space separately from the
values of its random variables. All proofs use elementary set operations,
real arithmetic and finite sums. Infinite sample spaces, integration,
limiting theorems and conditioning on null events require further results.

\section{Outcomes, events and probability masses}
\label{sec:p15-finite-probability-space}

\begin{hbdefinition}{def:p15-finite-probability-space}{Finite probability space}
Let \(\Omega\) be a finite nonempty set, whose elements \(\omega\) are
\emph{outcomes}. An \emph{event} is a subset of \(\Omega\); the set of
events is its power set \(2^\Omega\). A \emph{probability mass function}
on \(\Omega\) is a function \(p:\Omega\to[0,1]\) satisfying
\[
 \sum_{\omega\in\Omega}p(\omega)=1.
\]
It defines the \emph{probability measure}
\(\mathbb P:2^\Omega\to[0,1]\) by
\[
 \mathbb P(A)\defeq\sum_{\omega\in A}p(\omega).
\]
The triple \((\Omega,2^\Omega,\mathbb P)\) is the corresponding
\emph{finite probability space}. An empty sum is zero and an empty product
is one. In this chapter \(A^c\defeq\Omega\setminus A\) is complement
relative to the declared outcome space. Events \(A,B\) are
\emph{disjoint} when \(A\cap B=\varnothing\); a family is pairwise
disjoint when every two events with distinct indices are disjoint.

The \emph{support} of \(p\) is
\(\{\omega\in\Omega:p(\omega)>0\}\). A \emph{null event} has
probability zero. A property holds \emph{almost surely} when its event has
probability one. Outcomes outside the support may remain in \(\Omega\).
\index{probability!finite space}\index{probability!mass function}
\end{hbdefinition}

The source calls its atom-weight function a distribution function. Here
\emph{mass function} names those weights; it does not mean a density or a
cumulative distribution function. Uniform weights \(p(\omega)=1/|\Omega|\)
give \(\mathbb P(A)=|A|/|\Omega|\). Counting favorable outcomes uses this
uniformity assumption; a finite space alone does not supply it.

\begin{hbtheorem}{thm:p15-finite-probability-laws}{proved=annex:p15.finite-laws}{Finite probability identities}
For the space above:
\begin{enumerate}
\item \(\mathbb P(\varnothing)=0\), \(\mathbb P(\Omega)=1\), and
      \(0\leq\mathbb P(A)\leq1\).
\item If \(A\subseteq B\), then \(\mathbb P(A)\leq\mathbb P(B)\).
      If \(A_1,\ldots,A_m\) are pairwise disjoint, then
      \(\mathbb P(\bigcup_{i=1}^m A_i)=\sum_{i=1}^m\mathbb P(A_i)\),
      also for \(m=0\).
\item \(\mathbb P(A^c)=1-\mathbb P(A)\) and
      \(\mathbb P(A\cup B)=\mathbb P(A)+\mathbb P(B)-\mathbb P(A\cap B)\).
\item \(\mathbb P(A)=0\) exactly when every outcome in \(A\) has
      mass zero. Consequently \(\mathbb P(A)=1\) exactly when \(A\)
      contains the support.
\end{enumerate}
Conversely, let \(Q:2^\Omega\to[0,\infty)\) satisfy \(Q(\Omega)=1\)
and \(Q(A\cup B)=Q(A)+Q(B)\) for disjoint events. Then exactly one
mass function induces \(Q\), namely \(p(\omega)=Q(\{\omega\})\).
\end{hbtheorem}

\begin{hbdefinition}{def:p15-finite-random-variable}{Finite random variables and their laws}
A \emph{random variable} on this space with values in a finite set \(S\)
is a function \(X:\Omega\to S\). For \(D\subseteq S\), write
\[
 \{X\in D\}\defeq\{\omega\in\Omega:X(\omega)\in D\},\qquad
 \{X=x\}\defeq\{\omega\in\Omega:X(\omega)=x\}.
\]
These are events; \(\mathbb P(X=x)\) abbreviates
\(\mathbb P(\{X=x\})\). The \emph{law} of \(X\) is the mass
function \(p_X(x)\defeq\mathbb P(X=x)\) on \(S\).

For \(X:\Omega\to S\) and \(Y:\Omega\to T\), the \emph{joint
variable} is \((X,Y):\Omega\to S\times T\), and its \emph{joint law}
is \(p_{X,Y}(x,y)\defeq\mathbb P(\{X=x\}\cap\{Y=y\})\).
The individual laws \(p_X,p_Y\) are its \emph{marginal laws}.
For a function \(f\) defined on \(S\), \(f(X)\) denotes the
composition \(f\circ X\).
The chosen value sets may contain values with zero mass; the joint law may
vanish on pairs that the functions cannot attain.
\index{random variable!finite}\index{probability!joint and marginal laws}
\end{hbdefinition}

\begin{hbtheorem}{thm:p15-finite-law-projection}{proved=annex:p15.finite-law-projection}{Computing laws by grouping outcomes}
The preceding laws are probability mass functions, and
\[
\begin{aligned}
\mathbb P(X\in D)&=\sum_{x\in D}p_X(x),\\
p_X(x)&=\sum_{y\in T}p_{X,Y}(x,y),\qquad
p_Y(y)=\sum_{x\in S}p_{X,Y}(x,y).
\end{aligned}
\]
More generally, if \(f:S\to U\) is a function to a finite set, the law
of \(f(X)\) is
\[
 p_{f(X)}(u)=\sum_{\substack{x\in S\\f(x)=u}}p_X(x).
\]
\end{hbtheorem}

\section{Conditioning and inversion of evidence}
\label{sec:p15-finite-conditioning}

\begin{hbdefinition}{def:p15-finite-conditioning}{Conditioning on an event of positive probability}
For an event \(B\) with \(\mathbb P(B)>0\), define
\[
 \mathbb P(A\mid B)\defeq\frac{\mathbb P(A\cap B)}{\mathbb P(B)}.
\]
The conditional probability measure \(\mathbb P_B\) on the same event
set is \(\mathbb P_B(A)\defeq\mathbb P(A\mid B)\). Its mass function is
\[
 p_B(\omega)\defeq
 \begin{cases}
 p(\omega)/\mathbb P(B),&\omega\in B,\\
 0,&\omega\notin B.
 \end{cases}
\]
Thus outcomes outside \(B\) receive zero mass, and the relative weights
within \(B\) are retained. This is a defined update of a specified model
when the information is that \(B\) occurred; it is not a rule for an
unspecified change in the process generating the outcomes.
\index{conditional probability}
\end{hbdefinition}

\begin{hbtheorem}{thm:p15-finite-conditioning-laws}{proved=annex:p15.finite-conditioning}{Conditional measure, multiplication and successive conditioning}
The function \(p_B\) is a probability mass function with induced measure
\(\mathbb P_B\). In particular \(\mathbb P(B\mid B)=1\), and all
finite probability identities hold with \(\mathbb P_B\) in place of
\(\mathbb P\). Also
\[
 \mathbb P(A\cap B)=\mathbb P(A\mid B)\mathbb P(B).
\]
For events \(A_1,\ldots,A_n\), with \(n\geq1\), put
\(B_0=\Omega\) and \(B_i=\bigcap_{j=1}^i A_j\). If
\(\mathbb P(B_i)>0\) for \(0\leq i<n\), then the chain rule is
\[
 \mathbb P(B_n)=\prod_{i=1}^n\mathbb P(A_i\mid B_{i-1}).
\]
No independence is assumed. If \(\mathbb P(B\cap C)>0\), then
conditioning \(\mathbb P_B\) on \(C\) gives
\[
 \mathbb P_B(A\mid C)=\mathbb P(A\mid B\cap C).
\]
\end{hbtheorem}

\begin{hbwarning}{def:p15-null-conditioning}{A null conditioning event}
Take \(\Omega=\{a,b,c\}\) with masses \(1,0,0\), and let
\(B=\{b,c\}\). Although \(B\) is nonempty, \(\mathbb P(B)=0\).
The proposed expression \(\mathbb P(\{b\}\mid B)\) is the undefined
ratio \(0/0\). The multiplication equation would read \(0=q\cdot0\)
for every \(q\in[0,1]\), so it cannot determine a conditional value.
No convention assigning zero or one to that ratio is used here. The same
space also shows that an almost-sure event need not contain every outcome:
\(\{a\}\) has probability one but omits \(b,c\).
\end{hbwarning}

\begin{hbdefinition}{def:p15-finite-hypothesis-partition}{Hypotheses, likelihoods and posterior weights}
A \emph{finite partition} \(H_1,\ldots,H_m\) of \(\Omega\), with
\(m\geq1\), consists of nonempty, pairwise disjoint events whose union
is \(\Omega\). Regard these events as mutually exclusive and exhaustive
\emph{hypotheses}, and let \(E\) be an \emph{evidence event}.
The \emph{prior weights} are \(\pi_i\defeq\mathbb P(H_i)\). Put
\(J\defeq\{i:1\leq i\leq m,\ \pi_i>0\}\). For \(i\in J\),
the \emph{likelihood} of the evidence under \(H_i\) is
\(\ell_i\defeq\mathbb P(E\mid H_i)\). If \(\mathbb P(E)>0\),
the \emph{posterior weights} are \(\mathbb P(H_i\mid E)\).
Likelihoods are conditional probabilities of evidence; they need not sum
to one across the hypotheses.
For an event \(G\) with \(\mathbb P(G)<1\), its \emph{odds}
against \(G^c\) are \(\mathbb P(G)/(1-\mathbb P(G))\).
\index{Bayes' formula}\index{likelihood}\index{probability!prior and posterior}
\end{hbdefinition}

\begin{hbtheorem}{thm:p15-finite-bayes}{proved=annex:p15.finite-bayes}{Total probability and Bayes' formula}
For the partition and evidence above,
\[
 \mathbb P(E)=\sum_{i=1}^m\mathbb P(E\cap H_i)
             =\sum_{i\in J}\pi_i\ell_i.
\]
If \(\mathbb P(E)>0\), then
\[
 \mathbb P(H_i\mid E)=
 \begin{cases}
 \displaystyle\frac{\pi_i\ell_i}{\sum_{j\in J}\pi_j\ell_j},&i\in J,\\[1ex]
 0,&i\notin J.
 \end{cases}
\]
These posterior weights sum to one. A prior-null hypothesis remains null
after conditioning on evidence of positive probability.

For two hypotheses \(H,H^c\), put \(\pi\defeq\mathbb P(H)\), assume
\(0<\pi<1\), and write \(q_1\defeq\mathbb P(E\mid H)\) and
\(q_0\defeq\mathbb P(E\mid H^c)\). If \(q_0>0\), the
\emph{likelihood ratio} \(q_1/q_0\) and both odds below are defined, and
\[
 \frac{\mathbb P(H\mid E)}{\mathbb P(H^c\mid E)}
 =\frac{\pi}{1-\pi}\,\frac{q_1}{q_0}.
\]
Thus posterior odds equal prior odds times the likelihood ratio.
\end{hbtheorem}

\section{Independence and its limits}
\label{sec:p15-finite-independence}

\begin{hbdefinition}{def:p15-finite-independence}{Event and random-variable independence}
Events \(A,B\) are \emph{independent} when
\(\mathbb P(A\cap B)=\mathbb P(A)\mathbb P(B)\).
Events \(A_1,\ldots,A_n\) are \emph{mutually independent} when, for
every nonempty \(I\subseteq\{1,\ldots,n\}\),
\[
 \mathbb P\left(\bigcap_{i\in I}A_i\right)
       =\prod_{i\in I}\mathbb P(A_i).
\]
They are \emph{pairwise independent} when this holds for every two-element
\(I\). The empty-intersection convention is \(\bigcap\varnothing=\Omega\).

Finite-valued random variables \(X:\Omega\to S\) and
\(Y:\Omega\to T\) are \emph{independent} when
\[
 p_{X,Y}(x,y)=p_X(x)p_Y(y)\quad(x\in S,\ y\in T).
\]
Given an event \(H\) of positive probability, events or random variables
are \emph{conditionally independent given \(H\)} when the corresponding
definition holds for the conditional measure \(\mathbb P_H\).
\index{independence!pairwise and mutual}\index{conditional independence}
\end{hbdefinition}

\begin{hbtheorem}{thm:p15-finite-independence-laws}{proved=annex:p15.finite-independence}{Independence criteria and product construction}
\par\nobreak
If \(\mathbb P(B)>0\), events \(A,B\) are independent exactly when
\(\mathbb P(A\mid B)=\mathbb P(A)\). Every null event is independent
of every event. Independence of \(A,B\) implies independence of each
of the pairs \((A^c,B)\), \((A,B^c)\), and \((A^c,B^c)\).

The random variables \(X,Y\) are independent exactly when, for every
\(D\subseteq S\) and \(F\subseteq T\), the events \(\{X\in D\}\)
and \(\{Y\in F\}\) are independent.
For any mass functions \(a:S\to[0,1]\) and \(b:T\to[0,1]\) on
finite nonempty sets, the assignment \(p(x,y)\defeq a(x)b(y)\) is a mass
function on \(S\times T\); its coordinate variables are independent
with laws \(a,b\).
\end{hbtheorem}

\begin{hbexample}{def:p15-pairwise-not-mutual}{Pairwise independence and selection}
\par\nobreak
Let \(\Omega=\{(0,0),(0,1),(1,0),(1,1)\}\), with mass \(1/4\)
at each point. Let \(A\) say that the first coordinate is one,
\(B\) that the second is one, and \(C\) that the coordinates agree.
Each event has probability \(1/2\), and each pairwise intersection is
\(\{(1,1)\}\), with probability \(1/4\). Thus all three pairs are
independent. However,
\[
 \mathbb P(A\cap B\cap C)=\tfrac14
 \ne\tfrac18=\mathbb P(A)\mathbb P(B)\mathbb P(C),
\]
so the events are not mutually independent. Moreover,
\[
 \mathbb P(A\mid C)=\mathbb P(B\mid C)=\tfrac12,
 \qquad \mathbb P(A\cap B\mid C)=\tfrac12\ne\tfrac14.
\]
Selecting the equal-coordinate outcomes destroys the independence of
\(A,B\). Conditional independence must be checked in the specified
conditional measure.
\end{hbexample}

\begin{hbexample}{def:p15-alert-bayes-case}{A rare defect, an alert and repeated evidence}
\par\nobreak
Consider a stipulated finite model of a component and an inspection alert.
The four outcomes are defect/alert, defect/no alert, no defect/alert and
no defect/no alert. Assign their masses by dividing the following weights
by \(10000\). These are hypothetical model weights, not measured data.
\[
\begin{array}{c|rr|r}
 &E\ (\text{alert})&E^c\ (\text{no alert})&\text{total}\\\hline
 H\ (\text{defect})&90&10&100\\
 H^c\ (\text{no defect})&495&9405&9900\\\hline
 \text{total}&585&9415&10000
\end{array}
\]
The model gives
\[
 \mathbb P(H)=\tfrac1{100},\qquad
 \mathbb P(E\mid H)=\tfrac9{10},\qquad
 \mathbb P(E\mid H^c)=\tfrac1{20}.
\]
Total probability gives \(\mathbb P(E)=585/10000=117/2000\).
After an alert,
\[
 \mathbb P(H\mid E)=\frac{90}{585}=\frac2{13},\qquad
 \mathbb P(H^c\mid E)=\frac{11}{13}.
\]
The likelihood ratio is \(18\), the prior odds are \(1/99\), and the
posterior odds are \(18/99=2/11\), agreeing with these probabilities.
The likelihood \(9/10\) is not the posterior \(2/13\). If the prior
defect probability were instead \(1/10\), while the same conditional
alert probabilities held, the posterior would be
\[
 \frac{(1/10)(9/10)}{(1/10)(9/10)+(9/10)(1/20)}=\frac23.
\]
The different base rate changes the conclusion despite the same likelihoods.

For two alerts, extend the outcome space to triples
\((h,u,v)\in\{0,1\}^3\), where \(h=1\) means a defect and
\(u=1,v=1\) mean the first and second alerts. In this model let
\(H=\{h=1\}\), \(E_1=\{u=1\}\), and \(E_2=\{v=1\}\).
Set \(\pi_1=1/100\), \(\pi_0=99/100\),
\(r_1(1)=9/10\), \(r_0(1)=1/20\), and
\(r_h(0)=1-r_h(1)\). Assign mass
\[
 p_2(h,u,v)\defeq\pi_h r_h(u)r_h(v).
\]
Each \(r_h\) sums to one, so the total mass is
\(\sum_h\pi_h(\sum_u r_h(u))(\sum_v r_h(v))=1\).
Conditioning on \(h\) gives its two-coordinate product law. Thus this
explicit model assumes that \(E_1,E_2\) are conditionally independent
under each of \(H,H^c\), with the one-alert rates above. Two alerts give
\[
 \mathbb P(H\mid E_1\cap E_2)
 =\frac{(1/100)(9/10)^2}
        {(1/100)(9/10)^2+(99/100)(1/20)^2}
 =\frac{36}{47}.
\]
If the second reported alert is merely a copy of the first, so that
\(E_2=E_1\), both single-alert conditional probabilities still have the
same values, but the posterior after both reports is only \(2/13\).
Squaring likelihoods therefore requires the stated conditional independence;
the one-alert rates alone cannot justify it. No independence of the two
alerts in the unconditional population is asserted.
\end{hbexample}

\section{Finite expectation}
\label{sec:p15-finite-expectation}

\begin{hbdefinition}{def:p15-finite-expectation}{Expectation and event indicators}
A real-valued random variable on this finite space is a function
\(X:\Omega\to\mathbb R\). Its range is finite. Its
\emph{expectation}, or \emph{mean}, is
\[
 \mathbb E[X]\defeq\sum_{\omega\in\Omega}X(\omega)p(\omega).
\]
The \emph{indicator} of an event \(A\) is the random variable
\(\mathbf1_A(\omega)=1\) for \(\omega\in A\) and zero otherwise.
Arithmetic of random variables is pointwise: for example,
\((X+Y)(\omega)=X(\omega)+Y(\omega)\).
Likewise \(\{X=Y\}\defeq\{\omega:X(\omega)=Y(\omega)\}\) is an
event. Regard a real-valued variable as finite-valued on its range when
applying the preceding definitions of law and independence.
For a finite-valued \(X:\Omega\to S\) and a real-valued function
\(f:S\to\mathbb R\), \(f(X)\) denotes the composition \(f\circ X\).
\index{expectation!finite}\index{indicator of an event}
\end{hbdefinition}

\begin{hbtheorem}{thm:p15-finite-expectation-laws}{proved=annex:p15.finite-expectation}{Grouping, linearity and independent products}
\par\nobreak
For every finite-valued \(X\) and real-valued \(f\) as above,
\[
 \mathbb E[f(X)]=\sum_{x\in S}f(x)p_X(x).
\]
In particular \(\mathbb E[\mathbf1_A]=\mathbb P(A)\).
For real-valued \(X_1,\ldots,X_n\), real constants
\(c_0,c_1,\ldots,c_n\), and \(n\geq0\),
\[
 \mathbb E\left[c_0+\sum_{i=1}^n c_iX_i\right]
       =c_0+\sum_{i=1}^n c_i\mathbb E[X_i].
\]
This requires no independence. If \(X(\omega)\leq Y(\omega)\) for
every \(\omega\), then \(\mathbb E[X]\leq\mathbb E[Y]\).
If real-valued \(X,Y\) are independent, then
\(\mathbb E[XY]=\mathbb E[X]\mathbb E[Y]\).
\end{hbtheorem}

\begin{hbexample}{def:p15-expectation-with-dependence}{Counting and multiplying indicators}
For any events \(A_1,\ldots,A_n\), the random variable
\(N=\sum_{i=1}^n\mathbf1_{A_i}\) counts how many occur. Thus
\(\mathbb E[N]=\sum_i\mathbb P(A_i)\), even when the events overlap
or depend on one another. For instance, if \(A\) has probability
\(1/2\), then \(\mathbb E[\mathbf1_A+\mathbf1_A]=1\), although
the two summands coincide. In contrast,
\[
 \mathbb E[\mathbf1_A\mathbf1_A]=\mathbb E[\mathbf1_A]=\tfrac12
 \ne\tfrac14=\mathbb E[\mathbf1_A]^2.
\]
Linearity permits dependent summands; factorization of products needs its
own hypothesis.
\end{hbexample}

\begin{hbdefinition}{def:p15-finite-conditional-expectation}{Finite conditional expectation}
For a real-valued \(X\) and event \(B\) with \(\mathbb P(B)>0\),
\[
 \mathbb E[X\mid B]
 \defeq\sum_{\omega\in\Omega}X(\omega)p_B(\omega)
 =\frac{\sum_{\omega\in B}X(\omega)p(\omega)}{\mathbb P(B)}.
\]
This is a number. For a finite-valued \(Y:\Omega\to T\), put
\[
 g(y)\defeq
 \begin{cases}
 \mathbb E[X\mid\{Y=y\}],&p_Y(y)>0,\\
 0,&p_Y(y)=0.
 \end{cases}
\]
The random variable \(g(Y)\) is denoted \(\mathbb E[X\mid Y]\).
The displayed zero is a chosen value on null fibers; it is not a
definition of conditioning on those fibers. Changing it to another finite
real value changes the resulting variable only on a null event.
\index{conditional expectation!finite}
\end{hbdefinition}

\begin{hbtheorem}{thm:p15-finite-total-expectation}{proved=annex:p15.finite-total-expectation}{Averaging conditional means}
Let \(X:\Omega\to\mathbb R\) and \(Y:\Omega\to T\), with \(T\)
finite. For a partition \(H_1,\ldots,H_m\), put
\(J=\{i:\mathbb P(H_i)>0\}\). Then
\[
 \mathbb E[X]=\sum_{i\in J}\mathbb P(H_i)\mathbb E[X\mid H_i].
\]
For \(Z=\mathbb E[X\mid Y]\) and every \(D\subseteq T\),
\[
 \sum_{\omega:Y(\omega)\in D} Z(\omega)p(\omega)
 =\sum_{\omega:Y(\omega)\in D} X(\omega)p(\omega).
\]
In particular \(\mathbb E[\mathbb E[X\mid Y]]=\mathbb E[X]\).
Among variables of the form \(h(Y)\), with \(h:T\to\mathbb R\),
this averaging property determines \(Z\) almost surely.
\end{hbtheorem}

For the one-alert model in Example~\ref{def:p15-alert-bayes-case}, take
\(X=\mathbf1_H\). Then
\(\mathbb E[X\mid E]=2/13\), while
\(\mathbb E[X\mid E^c]=10/9415\). Averaging by the two alert outcomes
recovers the prior:
\[
 \frac{585}{10000}\frac{90}{585}
 +\frac{9415}{10000}\frac{10}{9415}
 =\frac{100}{10000}=\mathbb E[X].
\]
The conditional means vary with the evidence, but their probability-weighted
average equals the original mean.

\begin{hbwarning}{def:p15-finite-probability-limits}{Scope and use of a finite model}
Every probability and expectation above follows from the stipulated weights.
Using those numbers for an actual population additionally requires a justified
model of its outcomes, sampling and observations. The identities do not
estimate those weights, establish a causal effect, select a prior or prove
that an inspection process satisfies an independence assumption.

Finite real-valued expectations always exist. Their interpretation as limits
of empirical averages requires a separate law of large numbers with its own
hypotheses. This chapter provides neither continuous densities nor general
measurable spaces, integration, infinite product constructions, conditional
distributions on null events, convergence theorems or general conditional
expectation. Those foundations and their boundary cases remain separate
treatments.
\end{hbwarning}
